\section{Discussão}

Os resultados apoiam parcialmente a hipótese de que reduzir granularidade melhora a detecção. Material apresentou médias superiores a Fine, sobretudo em $F_1$, mas o IC pareado de AP inclui zero. Portanto, ``menos classes é melhor'' seria uma conclusão excessiva. Binary é claramente mais fácil, porém responde a outra pergunta operacional e perde toda informação sobre composição do resíduo.

A evidência mais forte vem da avaliação hierárquica. Quando uma caixa Fine recebe crédito após remapeamento, nenhuma propriedade geométrica ou confiança muda. O salto significativo nos dois níveis indica que muitas detecções contêm informação espacial útil, mas falham sob rótulos detalhados. Ainda assim, o AP binário controlado permanece abaixo de 0,5, e a decomposição registra muitas omissões; localização também está longe de resolvida.

A cauda longa oferece uma explicação complementar. As 28 classes raras têm, em média, somente 5,1 instâncias de treino e desempenho muito inferior às dez frequentes. Isso é consistente com a literatura de detecção de grande vocabulário \cite{gupta2019lvis}, mas aqui não demonstra causalidade isolada: tamanho, oclusão, contexto e semelhança visual também variam entre classes. A macroclasse Material compartilha dados entre rótulos relacionados e preserva uma informação potencialmente útil à triagem, constituindo compromisso plausível entre especificidade e robustez.

Há quatro ameaças principais à validade. Primeiro, uma única partição fixa não mede variação de amostragem; as três sementes medem somente a instabilidade de otimização. Segundo, o bootstrap cobre a semente 42 e condiciona-se ao teste observado. Terceiro, o mapeamento por material é uma decisão de projeto: compostos foram enviados a Other, e Organic contém apenas \textit{Food waste}. Quarto, não comparamos diretamente nossos números aos de trabalhos anteriores devido a diferenças de tarefa e protocolo. Estudos futuros devem repetir partições, avaliar agrupamentos alternativos e testar técnicas de balanceamento sem alterar simultaneamente outros componentes.
